\scanpage{43}
% archive sheet number: 42
déformation au\deleted{x} \deleted{\unsure{r}}cas du couple $(S_{\unsure{o|0}},S)$, \deleted{en} notant que l'idéal $p\underline{O}_S$ est bien un idéal à puissances divisées. On notera d'ailleurs que le théorème de déformations dans le cas \deleted{envisagé} général $(S,S')$ e\deleted{\unsure{n}}nvisagé dans son énoncé \deleted{\unsure{e}}\unsure{est|ést} en fait équivalent au cas particulier d'un couple $(S_o,S)$, comme il est formel v\deleted{\unsure{l}}ia $\mathrm{Crisloclib}(S)\handwritten{\xrightarrow{\approx}}\mathrm{Crisloclib}(S_o)$.

\inserted{\handwritten{\uline{Problème} 3\quad Trouver variante du th. de déformation pr les $p$-Groupes}}
% slanted
\margin{\handwritten{finis loc. libres (j'ai un candidat).}}

c) F-V-\uline{cristaux admissibles en car}. $p$. Revenons au cas d'un schéma de car. $p$ . Le cristal de Dieudonné $D^{\handwritten{*}}(G_o)$ étant muni d'une filtration canonique (de $D^{\handwritten{*}}(G_o)_{S_o}$), on pourrait craindre que le problème 1 est mal posé, et qu'il faudrait envisager le foncteur comme étant à valeurs dans des cristaux de Dieudonné filtrés. En fait, dans le cas envisagé ici, il se trouve que la filtration est uniquement déterminée en \deleted{y}termes de la structure de cristal de Dieudonné, comme on va l'expliquer maintenant.

Considérons le morphisme de topos annelés (qui aura été défini déjà précédemment, dans le par. d'exemples de cristaux, pour les relations entre cristaux et vecteurs de Witt)
\diagram{\begin{tikzcd}[ampersand replacement=\&]
\phi:\quad S_{o\,\mathrm{c}\deleted{\unsure{r}}\mathrm{ris}/\underline{F}_p}
 \arrow[r] \& S_o \arrow[l,bend right=38,"\handwritten{\varepsilon}"']\quad.
\end{tikzcd}}
Soit $M=D^{\handwritten{*}}(G_o)$, et soit $M_o$ sa restriction au site cristallin rel. à $\underline{F}_p$. Il se trouve que l\deleted{a\unsure{s|{}} suite} d'après $FV=p$, $VF=p$, les composé\deleted{\unsure{z}}s dans la suite
\diagram{\begin{tikzcd}[ampersand replacement=\&]
\deleted{\unsure{O|M}}M_o^{(p)}\arrow[r,"\handwritten{F_{M_o}}"] \&
M_o\arrow[r,"\handwritten{V_{M_o}}"] \&
M_o^{(p)}\arrow[r,"\handwritten{F_{M_o}}"] \& M_o
\end{tikzcd}}
sont nuls; en fait, on prouvera mieux, à savoir que \inserted{\handwritten{pour tout cristal de Dieud. $M$,}} la suite précédente est exacte (terme à terme sur tout épaississement à p.div. de car.$p$ de $S_o$) et on obtient ainsi des sous-cristaux en modules loc. libres
\[
\operatorname{Ker} V_{M_o}=\operatorname{Im} F_{M_o}\handwritten{\subset}M_o,\quad
\operatorname{Ker} F_{M_o}=\operatorname{Im} V_{M_o}\handwritten{\subset}M_o^{(p)}\quad .
\]
Ceci dit, \deleted{\unsure{le|de} sous-Module $\mathrm{Fil}^1$} notons que $M_{S_o}=\handwritten{\varepsilon^*}(M_o)$, donc
\[
\phi^{\handwritten{*}}(M_{S_o})\handwritten{\simeq} M_o^{(p)}
\]
donc on a
\[
\phi^{\handwritten{*}}(\deleted{\unsure{E|F}}\mathrm{Fil}^1(M_{S_o}))\handwritten{\subset} M_o^{(p)}\quad .
\]
Ceci posé, on montre qu'on a
